\documentclass[a4paper,12pt]{article}

% --- Wichtige Pakete ---
\usepackage[utf8]{inputenc}       % Erlaubt Umlaute (ä, ö, ü)
\usepackage[ngerman]{babel}       % Deutsche Trennung und Begriffe (z.B. "Datum")
\usepackage[T1]{fontenc}          % Schöne Schriftarten-Kodierung
\usepackage{amsmath,amssymb,amsfonts} % Für mathematische Formeln
\usepackage{graphicx}             % Für das Einfügen von Bildern
\usepackage{geometry}             % Für die Seitenränder
\usepackage{fancyhdr}             % Für schöne Kopf- und Fußzeilen

% --- Fehlerbehebungen ---
\setlength{\headheight}{14.5pt}   % Behebt den fancyhdr-Fehler
\DeclareUnicodeCharacter{2212}{\ensuremath{-}}    	% Behebt Minuszeichen-Fehler
\DeclareUnicodeCharacter{2219}{\ensuremath{\cdot}} 	% Behebt Malzeichen-Fehler

% --- Seitenränder einstellen ---
\geometry{left=2.5cm, right=2.5cm, top=3cm, bottom=3cm}

% --- Kopf- und Fußzeile einrichten ---
\pagestyle{fancy}
\fancyhf{}
\lhead{Name: \textbf{Fabian Seemann}}       	 % Dein Name links oben
\chead{Studiengang: \textbf{BiF}}                % Deine Klasse in der Mitte
\rhead{\today}                              	 % Aktuelles Datum rechts oben
\lfoot{Fach: ALGOS}                    	 	% Fach links unten
\rfoot{Seite \thepage}                      % Seitenzahl rechts unten
\renewcommand{\headrulewidth}{0.4pt}        % Linie unter der Kopfzeile
\renewcommand{\footrulewidth}{0.4pt}        % Linie über der Fußzeile

\begin{document}
	
	% --- Titel des Arbeitsblattes ---
	\begin{center}
		\LARGE{\textbf{Lösungen zum Arbeitsblatt 1:}} \\[0.2cm]
		\large{Thema: Komplexität, Suche, Hashing}
	\end{center}
	\vspace{0.5cm}
	
	% --- Aufgabe 1 ---
	\section*{Aufgabe 1.1: O-Notation Rechenregeln}
	Bestimmen/Vereinfachen Sie die Laufzeitkomplexität der folgenden Ausdrücke:
	
	\noindent
	\begin{tabular}{llp{9cm}}
		\textbf{Gleichung:} & $O(10n) = O(n)$ & \\
		\textbf{Regel:}     & Konstantenregel & (Skalierung) \\
		\textbf{Erklärung:} & \multicolumn{2}{p{11cm}}{Der Faktor 10 ist eine feste Konstante. Sie verändert sich nicht, wenn $n$ wächst. Uns interessiert nur die Art des Wachstums (linear). Die genaue Steilheit der Kurve ist in der O-Notation egal. Daher fällt die 10 weg.}
	\end{tabular}
	
	
	% --- Aufgabe 2 ---
	\section*{Aufgabe 1.2: Pseudocode}
	
	
	
	% --- Aufgabe 3 ---
	\section*{Aufgabe 1.3: Typische Komplexität}
	
	
\end{document}
